\begin{obereflection}
\textbf{Latihan 1.1}
1. Jelaskan perbedaan antara \textit{confidentiality} dan \textit{integrity}
   dengan contoh pengiriman nilai ujian mahasiswa.
2. Mengapa Prinsip Kerckhoffs menyatakan bahwa algoritma harus publik dan hanya
   kunci yang harus rahasia? Apa risikonya jika algoritma juga dirahasiakan?
3. Sebuah perusahaan menyimpan basis data pelanggan dalam berkas yang telah
   di-\textit{encode} Base64. Manajemen mengklaim data itu ``sudah terenkripsi''.
   Bantah klaim tersebut.
4. Jelaskan empat layanan kriptografi dan berikan satu contoh penerapan untuk
   masing-masing.
5. Mengapa kriptografer masih langka di Indonesia padahal kebutuhan tenaga ahli
   bidang ini besar? Sebutkan dua faktor penghambat.
6. Bandingkan kriptografi simetri dan nir-simetri dari jumlah kunci yang harus
   dikelola. Hitung jumlah kunci untuk 100 pihak pada masing-masing sistem.
\end{obereflection}
